\chapter{INTRODUCTION} 

\section{Background of the Study}
The increasing demand for mental health support has significantly accelerated efforts to develop Large Language Models (LLMs) as accessible psychological counselors \cite{lee2024cactus}. Within therapeutic approaches, Cognitive Behavioral Therapy (CBT) stands out as one of the most effective non-pharmacological methods, proven to enhance well-being and reduce perceived stress \cite{shahrokhian2021effects}. To enable LLMs to function as effective CBT practitioners, it is crucial to train them on realistic, high-quality data that emulates real-world therapeutic interactions, a challenge addressed by recent efforts in dataset creation such as CACTUS \cite{lee2024cactus}.

\section{Statement of the Problem}
However, a critical shortcoming of current text-based LLMs is their inability to process the nuances of verbal communication; they excel at understanding what is said, but fail to grasp how it is said \cite{sun2024beyond}. This limitation ignores a vital channel of information, as clinicians and researchers have long recognized that non-verbal, paralinguistic vocal cues are powerful indicators of an individual's mental state \cite{cummins2015review}. The human voice is a highly sensitive output system, affected by the cognitive and physiological changes associated with conditions like depression, which manifest in features related to prosody, voice quality, and spectral characteristics \cite{cummins2015review}. Recent work in other domains has shown that integrating vocal cues such as pitch, loudness, and duration can significantly improve an LLM's decision-making capabilities, especially in situations with ambiguity \cite{sun2024beyond}. The critical gap, therefore, lies in the application of this principle to therapeutic AI; there is a lack of a framework that integrates real-time vocal cue analysis into an interactive AI to deliver in-the-moment therapeutic actions.

\section{Research Questions}
\begin{enumerate}
    \item \textbf{RQ1 -- Dissonance Advantage.} Does the Dissonance-Aware therapist framework, which computes emotional mismatch (VA delta) between text and speech, consistently outperform Baseline (text-only), Emotion (text+VA), and Multimodal (text+speech+VA, without dissonance) methods on measurable therapeutic quality metrics?

    \item \textbf{RQ2 -- Backbone Robustness.} Is the advantage of the Dissonance-Aware framework robust across different LLM backbones? Specifically, does the dissonance signal provide consistent performance gains when the underlying therapist LLM is changed (e.g., GPT-4o mini vs.\ Claude Sonnet 4.6 vs.\ Qwen-2.5-72B-Instruct vs.\ DeepSeek V3.2)?
\end{enumerate}

\section{Hypotheses}

\begin{enumerate}
    \item \textbf{H1 (Dissonance Advantage -- RQ1):} The Dissonance-Aware therapist will achieve the highest (or tied-highest) scores on all seven evaluation metrics, Understanding, Interpersonal Effectiveness, Affective Bond, CTRS Collaboration, Guided Discovery, Focus, and Strategy, compared to Baseline, Emotion (Text-Only), and Multimodal (Vocal-Aware) conditions.

    \item \textbf{H2 (Backbone Robustness -- RQ2):} The performance advantage of the Dissonance-Aware framework over the other three conditions will persist across different LLM backbones, with the delta-driven dissonance signal providing a consistent, architecture-independent improvement.
\end{enumerate}

\section{Objectives of the Study}

\begin{enumerate}
    \item To construct a multimodal emotional analysis pipeline that independently extracts Valence-Arousal from text (via vad-bert) and speech (via WavLM), computes emotional dissonance as the VA delta between modalities, and injects that signal into therapist LLM prompts alongside vocal prosodic descriptors.

    \item To implement and compare four therapist generation conditions across multiple LLM backbones, isolating whether the dissonance signal's therapeutic benefit is model-dependent or architecture-agnostic.
\end{enumerate}

\section{Methodology Overview}

To address these questions, this thesis proposes a multimodal dissonance-aware therapist framework. The framework extracts emotional signals from two parallel channels: text-derived Valence-Arousal from the client's words, and speech-derived Valence-Arousal from their vocal delivery. A dissonance score, computed as the delta between these two channels, flags moments where what the client says and how they say it diverge. Four therapist generation conditions, each providing a different subset of these emotional signals to the same underlying LLM, are compared across four LLM backbones on 100 shared-problem-seed CBT dialogues per cell (1,600 sessions) through a three-layer blind evaluation: absolute GPT-4o grading on seven therapeutic quality metrics, per-seed paired statistics, and swap-consistent pairwise preference comparisons with a cross-family second judge. This controlled comparison isolates the specific contribution of the dissonance signal to therapist response quality.

\section{Scope and Limitations}

\textbf{In Scope (What This Research Will Do):}
\begin{itemize}
    \item \textbf{Develop a Framework:} Propose and build a multimodal therapist framework that integrates text and speech emotional signals to detect and respond to emotional dissonance in CBT-style dialogue.
    \item \textbf{Synthetic Dataset:} Create and utilize a synthetically generated dataset of 100 multi-turn CBT counseling dialogues per condition and backbone (4 backbones $\times$ 4 conditions $\times$ 100 shared problem seeds), produced via CACTUS-style persona simulation with stateful client memory and Zonos TTS synthesis with state-controlled prosody.
    \item \textbf{Multimodal Condition Comparison:} Compare four therapist generation conditions, Baseline, Emotion (Text-Only), Multimodal (Vocal-Aware), and Dissonance-Aware, to quantify the contribution of each emotional signal type.
    \item \textbf{Automated Evaluation:} Assess therapist response quality using GPT-4o as a zero-shot evaluator across seven metrics spanning Therapist Skills, Client Alliance, and the Cognitive Therapy Rating Scale.
\end{itemize}

\textbf{Out of Scope (What This Research Will NOT Do):}
\begin{itemize}
    \item \textbf{Not a Clinical Tool:} This project is a computer science exploration. The resulting AI is not a medical device and is not intended to replace professional human therapists or provide clinical diagnoses.
    \item \textbf{No Real Patient Data:} For ethical and privacy reasons, this research will not use any real-world patient or clinical data. All training and testing will be performed on the synthetically generated dataset.
    \item \textbf{Limited to English:} The current scope is limited to processing and generating English-language speech based on the available datasets (CACTUS, Zonos).
    \item \textbf{No Other Modalities:} The analysis is strictly focused on vocal and textual information. We will not analyze other cues such as facial expressions or body language.
\end{itemize}

\section{Key Findings}

Across 1,600 evaluated sessions and 800 pairwise comparisons, three findings stand out. First, the Dissonance-Aware condition achieved the highest score on all seven therapeutic quality metrics on GPT-4o mini, with statistically significant gains over the text-only Baseline on the CTRS craft dimensions (Guided Discovery $+0.24$, Focus $+0.34$, Strategy $+0.17$; pairwise net win-rate $0.680$, $p = 0.0001$), and a significant advantage on DeepSeek V3.2 (pairwise $p = 0.00004$) that absolute scores alone could not detect. Second, the Multimodal (Vocal-Aware) condition --- the same vocal inputs without dissonance framing --- \emph{regressed below Baseline} on both weaker backbones, while Dissonance-Aware defeated it decisively (GPT-4o mini: pairwise net win-rate $0.800$, $p < 10^{-4}$): the dissonance computation, not the vocal data itself, is the active ingredient. Third, the advantage is headroom-gated: neutral on Qwen's already-strong text-only therapist and unmeasurable at Claude's scoring ceiling, with no backbone showing a significant loss for Dissonance-Aware anywhere. These findings support the central thesis claim with explicit boundary conditions: emotional dissonance, framed therapeutically, converts noisy multimodal signals into measurable improvements in CBT response quality precisely for the cost-constrained models most likely to be deployed.

\section{Contributions}

The primary contributions of this work are: (1) a multimodal framework that independently extracts and aligns text and speech emotional representations (VA) for therapist response generation; (2) the operationalization of emotional dissonance, computed as a VA delta between modalities, as a prompt-level therapeutic signal --- together with evidence that this framing, rather than the raw vocal metadata, produces the quality gain; (3) a controlled four-condition, four-backbone experimental design, paired on shared problem seeds with stateful simulated clients, that isolates the contribution of each emotional metadata type; and (4) a reproducible three-layer blind evaluation protocol --- absolute seven-metric grading, per-seed paired statistics, and swap-consistent pairwise preference with cross-family judge reliability measurement --- that remains discriminative when absolute LLM-judge scores saturate near ceiling. The remainder of this paper is organized as follows: Chapter~2 reviews related work, Chapter~3 details the methodology, Chapter~4 describes the experimental setup, and Chapter~5 presents the results, with Chapter~6 discussing implications and limitations.